% Copyright (c) 2026 Twin Vector Labs LLC.
% All rights reserved.
\documentclass[11pt]{article}

\usepackage[margin=1in]{geometry}
\usepackage{amsmath,amssymb,amsthm}
\usepackage{mathtools}
\usepackage{enumitem}
\usepackage{booktabs}
\usepackage{array}
\usepackage[colorlinks=true,linkcolor=blue,urlcolor=blue]{hyperref}

\newcommand{\C}{\mathbb{C}}
\newcommand{\Z}{\mathbb{Z}}
\newcommand{\rep}[1]{\mathbf{#1}}
\newcommand{\abs}[1]{\left\lvert #1 \right\rvert}
\newcommand{\norm}[1]{\left\lVert #1 \right\rVert}
\newcommand{\sigmacol}{\sigma}
\newcommand{\Pin}{P_{\mathrm{in}}}
\newcommand{\Pout}{P_{\mathrm{out}}}

\title{\bfseries The Proton So Far:\\
an emergent color singlet, its observables, and lattice convergence}
\author{Tessera --- cobordism subsystem (\#398 / \#400 / \#404, under the \#380 epic)}
\date{}

\begin{document}
\maketitle

\begin{abstract}
This is a status report on building a proton from quarks in the tessera
cobordism/Regge subsystem. Three results are in hand. (1) \textbf{The emergent color
singlet} (\#398): on the trivalent junction $W_{ABC}$ with \emph{symmetrically} placed
register windows, the input$\to$result transport intertwines the color $\Z_3$, and a
natural color-symmetric quark input transports to the net-color-zero singlet line
$\langle(1,\omega,\omega^2)\rangle$ with overlap $0.999$ --- the output's location is
forced by symmetry, not hand-placed. (2) \textbf{Emergent observables} (\#400), read
off the relaxed singlet: the color charge $\sigmacol\to 0$ (confinement), a length
scale, and a (proxy) mass, giving the dimensionless $r\!\cdot\!m \approx 6.9$ (the
physical proton's value is $\approx 4.0$). (3) \textbf{Lattice convergence} (\#404):
the granularity is tunable (a frequency-$N$ geodesic base); refining it shrinks the
intertwiner residual while the singlet overlap sharpens to $1$ --- consistent with the
small imperfections being discretization artifacts, not symmetry breaking. The honest
assessment: we have the \emph{color-confinement skeleton} of a baryon --- emergent,
structurally correct, with a scale-free observable in the right ballpark --- on a
coarse lattice; not yet a quantitatively faithful proton (the size is still
lattice-pinned, the mass is a curvature proxy, and the electroweak quantum numbers are
absent). All numbers are measured through the production C++ \texttt{TransportCobordism}.
\end{abstract}

\tableofcontents

\section{The program}
The proton-from-quarks program relaxes a Lorentzian Regge cobordism with pinned
boundary color states and reads the \emph{emergent} result off the relaxed geometry
--- nothing is imposed. The color singlet of $SU(3)$ is the alternating invariant
$\varepsilon_{ijk}$, which is absent from $\rep{3}\otimes\rep{3}$ and first appears in
$\rep{3}\otimes\rep{3}\otimes\rep{3}$: a baryon is irreducibly a three-quark object,
so pairwise (bipartite) merges cannot reach it (\#382). The resolution is the
trivalent junction $W_{ABC}$ (\#396): one connected geodesic icosahedron ($S^2$)
minus twelve vertex-disjoint hole triangles in four \emph{windows} of three --- the
quark inputs $A,B,C$ and the emergent result $R$ --- extruded $\times I$. Distinct
holes are independent cycles ($b_1=11$), so the inputs do not average, and charge is
conserved at the junction by the global Stokes relation
$\sigmacol_R = -(\sigmacol_A+\sigmacol_B+\sigmacol_C)$.

\section{The emergent color singlet (\#398)}
The four windows are placed \emph{symmetrically}: $A_4$, a tetrahedral subgroup of the
icosahedral rotation group, partitions one tetrahedral set of icosahedron vertices
into four blocks of three (the window corners), and permutes the four blocks
transitively --- so the windows are $A_4$-equivalent. Write $g$ for the $A_4$
three-cycle that fixes $R$ and cycles $A\to B\to C$; on the result holes it is the
cyclic color $\Z_3$ (the cyclic subgroup of the hole-permutation $S_3$/Weyl symmetry,
which is why only the cube roots of unity appear). Let $\Pin$ (on the $9$-dimensional
input space over $A,B,C$) and $\Pout$ (on the $3$-dimensional result space over $R$'s
holes) be its signed permutation action on hole periods, sharing an orientation
convention, and let $M:\,V_{\mathrm{in}}\to V_{\mathrm{out}}$ be the input$\to$result
transport. Then $M$ \textbf{intertwines} the symmetry, to the residual of Table~2,
\begin{equation}
  M\,\Pin \approx \Pout\,M .
  \label{eq:intertwine}
\end{equation}
$\Pout$ is the color $\Z_3$ (eigenvalues $\{1,\omega,\omega^2\}$); its
$\omega$-eigenline is the singlet direction $\langle(1,\omega,\omega^2)\rangle$. This
is the discrete \emph{net-color-zero} ($\sigmacol=0$) surrogate for the SU(3) singlet
--- it is the confinement condition (orthogonality to the charge mode $(1,1,1)$),
the $\Z_3$ shadow of $\varepsilon_{ijk}$, not literally $\varepsilon_{ijk}$ (which is
slot-permutation invariant). By \eqref{eq:intertwine}, the $\omega$-eigenline of
$\Pin$ --- the natural color-symmetric quark input --- is carried into the singlet
line up to the residual: the output's location is \emph{forced} (not hand-placed),
even though the input representation is the symmetry-distinguished choice.

Three ingredients are jointly necessary (symmetric windows alone are not enough): a
symmetry-respecting metric, the input in the $\omega$-representation of $g$, and a
$g$-coherent signed read-out. With them, the measured overlap is $0.999$ (the $0.001$
deficit is the intertwiner residual; a greedy, geometrically inequivalent placement
reaches only $\sim0.74$). Details are in the companion note \emph{The Proton Color
Singlet from Symmetric Register Windows}.

\section{Emergent observables (\#400)}
With the singlet in hand we read the proton's observables \emph{off the relaxed
geometry}, emergent-first, reusing the existing measurement methods. The singlet
survives relaxation ($0.995$), so all reads come from one relaxed geometry.

\begin{center}
\small
\setlength{\tabcolsep}{5pt}
\begin{tabular}{@{}lcl@{}}
\toprule
observable & value & how it is read \\
\midrule
color charge $\sigmacol$ & $0.06 \approx 0$ & projection onto $\Pout$'s $+1$ (charge) mode --- confinement\\
singlet component & $0.995$ & the $\omega$-mode (the proton) \\
radius $r$ (RMS edge / lattice scale) & $1.063$ & $\sqrt{\langle\ell^2\rangle_{\ell^2>0}}$ over relaxed spacelike edges\\
mass $B$ (curvature anchor) & $6.53$ & shell-summed Re-deficit (the \#352 method)\\
mass $A$ (alt.\ proxy) & $228.8$ & $\abs{\text{dual Regge action}}$\\
$r\!\cdot\!m$ (anchor $B$) & $\mathbf{6.94}$ & vs the physical proton's $\approx 4.0$\\
$r\!\cdot\!m$ (alt.\ proxy $A$) & $243$ & \\
\bottomrule
\end{tabular}
\end{center}

\noindent The color charge is the $g$-invariant ($+1$) mode of the result; the singlet
is the orthogonal $\omega$-mode, so $\sigmacol\to 0$ is confinement --- this realizes
the Stokes relation $\sigmacol_R=-(\sigmacol_A+\sigmacol_B+\sigmacol_C)$, a
color-symmetric input summing to zero result charge. \textbf{The two masses are
different, un-normalized proxies} --- a dual-action magnitude ($A$) versus a
shell-summed real deficit ($B$) --- and are not expected to agree; the $\sim$35$\times$
gap is a normalization difference, not a consistency check. $B$ is the established
\#352 curvature anchor used below; $A$ is reported for scale only. The dimensionless
$r\!\cdot\!m$ is the scale-free observable to compare across unit systems: the physical
proton has $r_p m_p \approx (0.84~\text{fm})(938~\text{MeV}/\hbar c) \approx 4.0$ ---
its radius is about four reduced Compton wavelengths. Our $6.9$ (anchored on the
curvature mass $B$) is the same order of magnitude, and a large improvement on the
prior crude $\sim 73$ from the bare-quark merge (\#352); but the comparison is
heuristic (proxies; see \S\ref{sec:assess}).

\section{Lattice convergence (\#404)}
The base granularity is tunable: a frequency-$N$ geodesic icosahedron
(\texttt{set\_frequency(N)}, $N=2$ the $42$-vertex default), with the same
$A_4$-orbit windows generated from the symmetry at any $N$. The small imperfection in
\eqref{eq:intertwine} is a fixed-resolution triangulation artifact (the $\times I$
extrusion is not $g$-symmetric even though the base surface and metric are), so it
should shrink with resolution while the singlet overlap sharpens toward~$1$:

\begin{center}
\setlength{\tabcolsep}{8pt}
\begin{tabular}{@{}cccc@{}}
\toprule
$N$ & cobordism vertices & $\norm{M\Pin-\Pout M}/\norm{M}$ & singlet overlap\\
\midrule
2 & 126 & $4.26\times10^{-2}$ & $0.9991$--$0.9998$\\
3 & 276 & $2.09\times10^{-2}$ & $0.9996$--$0.99999$\\
4 & 486 & $1.69\times10^{-2}$ & $0.9997$--$1.0000$\\
\bottomrule
\end{tabular}
\end{center}

\noindent The overlap is essentially exact at every $N$; the residual decreases
monotonically, but the $N{=}3\to4$ step is small ($\sim$19\%), so whether it reaches
zero or floors near $\sim$1.7$\times10^{-2}$ is not yet established. The strong
evidence that the imperfection is a discretization artifact, not symmetry breaking, is
the overlap $\to 1$. $N=2$ reproduces the \#398 base exactly (backward-compatible: the
junction suite passes $20/20$, and \texttt{set\_frequency} rejects $N<2$, the
$42$-vertex floor).

\section{Are we looking at a proton?}
\label{sec:assess}
A careful answer: we have the \textbf{color-confinement skeleton} of a baryon, not a
quantitatively faithful proton.

\paragraph{What is genuinely real.} The color singlet. Three colored quarks bind into
a color-neutral, confined state --- the defining property of a baryon --- and its
\emph{location} emerges from a symmetric input, forced by the geometry's $\Z_3$
symmetry, never hand-placed ($\sigmacol\to 0$, singlet $0.995$, robust through
relaxation and sharpening with resolution).

\paragraph{What is still a proxy.} The \emph{radius} is essentially the lattice scale
(the RMS edge length $r\approx 1.06$ on the uniform seed): the geometry did not develop
its own localized size, so $r$ carries little physical information yet, and
$r\!\cdot\!m$ is dominated by the mass term. The \emph{mass} is a curvature
(deficit-angle) proxy, not a calibrated rest mass. So the $6.9$ vs $4.0$ agreement is
an order-of-magnitude correspondence for a proxy, not a measurement. And this is only
the \emph{color} sector: the proton is also electric charge $+1$, spin-$\tfrac12$, and
flavor $uud$, none of which this yet captures; the lattice is also coarse.

\paragraph{Verdict.} A proton-shaped object: the right confinement and color-singlet
structure, arrived at emergently with the correct symmetry, and a scale-free
observable in the right ballpark --- the QCD skeleton of a proton on a coarse lattice,
not the full particle.

\section{Outlook}
Three follow-ons (the \#380 epic) would push it toward ``actual'':
\begin{itemize}[nosep]
  \item \textbf{\#403 --- let the size emerge.} Relax from a localizing, $g$-invariant
  seed (the van Raamsdonk / entanglement seed, available via
  \texttt{set\_entangled\_metric}) so a genuine bound-state length scale emerges and
  the radius is no longer lattice-pinned; then $r\!\cdot\!m$ becomes a clean probe.
  This composes with the finer lattice (\#404). (A reconnaissance confirms the
  entanglement seed develops a strong interior/exterior contrast where the uniform
  seed stays flat.)
  \item \textbf{\#405 --- electroweak quantum numbers.} Read the emergent electric
  charge $+1$ (the $uud$ charges $+\tfrac23,+\tfrac23,-\tfrac13$), then spin and flavor.
  \item \textbf{\#361 --- animation.} Visualize the relaxation: primal/dual geometry,
  curvature gradient, null (photon) edges, live residuals.
\end{itemize}

\end{document}
